\begin{abstract}
Large language models increasingly operate in multi-step settings where strategic misrepresentation can improve local outcomes, but most honesty evaluations remain single-turn and static. We introduce \textbf{BS-Bench}, a reproducible benchmark for deception, lie detection, and instruction compliance in multi-agent play by having four LLM agents play the card game \emph{Bullshit}. The environment is useful because lying is legal, the truth of each claim is objectively verifiable from hidden cards, and outcomes depend on both bluffing and challenge behavior under uncertainty. We evaluate six hosted models under four prompt conditions: a low-strategy control, deception allowed, asymmetric fairness, and an honesty mandate. The frozen comparable cohort contains 600 winner-terminated games with no mixed-cohort or turn-cap exclusions. Under the honesty mandate, mean model-level lie frequency falls from 39.6\% to 27.2\%. But challenge frequency also falls from 45.3\% to 21.8\%, while lie success rises from 11.0\% to 34.7\%. The asymmetric-fairness condition yields heterogeneous lie-frequency shifts from -7.4 to +3.3 points rather than a universal restraint effect. These results suggest that prompt framing changes the enforcement regime of the table, not just individual willingness to lie.
\end{abstract}
